% Source: https://artofproblemsolving.com/community/c3h1689240_question
% Stable identity: aops-topic-1689240-q01
\textbf{(a)} Let $s=x+\dfrac{1}{x}$. Then $s^2=x^2+2+\dfrac{1}{x^2}=5$, so $s=\pm\sqrt5$; both signs occur, since $x$ and $-x$ both satisfy the condition (for example $x=\dfrac{\sqrt5\pm1}{2}$ give $s=\sqrt5$). Then
\[
x^3+\dfrac{1}{x^3}=\left(x+\dfrac{1}{x}\right)\left(x^2-1+\dfrac{1}{x^2}\right)=s(3-1)=2s .
\]
So $x^3+\dfrac{1}{x^3}=\boxed{\pm2\sqrt5}$: it is $2\sqrt5$ for positive $x$ and $-2\sqrt5$ for negative $x$.

\textbf{(b)} Divide the numerator and denominator by $x^2$:
\[
\dfrac{x^2}{(x^2+1)^2}=\dfrac{1}{\dfrac{(x^2+1)^2}{x^2}}=\dfrac{1}{x^2+2+\dfrac{1}{x^2}}=\dfrac{1}{3+2}=\boxed{\dfrac{1}{5}}.
\]
